% Vietnamese translation for OI Public Library
% Source-PDF: IOI中国国家候选队论文/2000/徐静论文.pdf
\documentclass[11pt]{article}
\usepackage{oipl-vn}

\title{Xây dựng và chuyển hóa mô hình đồ thị}
\author{Xu Jing}
\date{}

\begin{document}
\maketitle

\origtitle{Construction and transformation of graph-theory models}
\sourcepath{IOI China candidate team papers / 2000 / Xu Jing.pdf}

\renewcommand{\contentsname}{Mục lục}
\tableofcontents

\paragraph{Từ khóa.}
Mô hình đồ thị; xây dựng; chuyển hóa.

\section*{Tóm tắt}

Bài viết này chủ yếu bàn về việc xây dựng và chuyển hóa mô hình đồ thị, gồm
bốn phần.

Phần thứ nhất là dẫn nhập, nêu vị trí của mô hình hóa đồ thị trong toàn bộ
kỳ thi tin học, so sánh điểm giống và khác giữa mô hình đồ thị với các mô
hình toán học khác, đồng thời chỉ ra sự cần thiết của việc nghiên cứu, tổng
kết tư tưởng, phương pháp và kỹ xảo mô hình hóa đồ thị.

Phần thứ hai nêu hai điểm then chốt trong xây dựng mô hình đồ thị: chọn bỏ
các yếu tố trong nguyên mẫu một cách thích hợp, và chọn hệ lý thuyết phù
hợp. Mỗi điểm đều được phân tích bằng ví dụ, rồi từ đó rút ra nguyên tắc
chung của mô hình hóa đồ thị: chính xác, rõ ràng, ngắn gọn.

Phần thứ ba thảo luận hai phương pháp được dùng khá rộng rãi trong chuyển hóa
mô hình đồ thị: chuyển hóa bằng tách và chuyển hóa bằng phần bù, trong đó tập
trung phân tích phương pháp thứ nhất. Chuyển hóa bằng tách được chia thành ba
loại: đỉnh thành cạnh, đỉnh thành đỉnh, và cạnh thành cạnh; loại thứ hai được
phân tích chi tiết hơn.

Phần thứ tư tổng kết toàn văn và nêu sự cần thiết của việc tiếp tục nghiên
cứu mô hình đồ thị.

\section{Dẫn nhập}

Thi tin học lấy giải bài làm trung tâm. Trong toàn bộ quá trình giải bài, một
bước quan trọng là mô hình hóa toán học; điều bài viết này bàn tới là một
nhánh của mô hình hóa toán học: mô hình hóa đồ thị.

Mô hình hóa đồ thị là quá trình trừu tượng hóa và đơn giản hóa một số sự vật
khách quan, rồi dùng đồ thị\footnote{Trong bài viết này, ``đồ thị'' chỉ hình
gồm một số đỉnh phân biệt và các cạnh nối một số cặp đỉnh trong đó; không bao
gồm các hình minh họa nói chung.} để mô tả đặc trưng và quan hệ nội tại của
chúng.

Mục đích của việc xây dựng mô hình đồ thị cũng giống các mô hình toán học
khác: đơn giản hóa vấn đề, làm nổi bật trọng điểm, để có thể nghiên cứu bản
chất vấn đề sâu hơn. Mục tiêu giải có thể là bài toán tối ưu, cũng có thể là
bài toán tồn tại hoặc bài toán xây dựng. Hơn nữa, giống mô hình hình học hay
mô hình vận trù học, trong quá trình xây dựng mô hình đồ thị ta cũng cần dùng
các khái niệm và công cụ toán học cơ bản như tập hợp, ánh xạ và hàm số.

Nhưng mô hình đồ thị cũng rất khác các mô hình khác ở phương pháp nghiên cứu.
Chẳng hạn, ta có thể dùng các thuật toán đồ thị điển hình để giải mô hình đồ
thị, hoặc dựa trên lý thuyết đồ thị cơ bản để phân tích tính chất của mô
hình. Những thuật toán và lý thuyết đặc thù ấy không có trong các mô hình
khác; hơn nữa, trong các mô hình khác cũng hiếm khi có một cấu trúc trực quan
như đồ thị để mô tả vấn đề.

Khi học lý thuyết đồ thị, thông thường ta học qua sách. Nhưng sách thường chỉ
giới thiệu các yếu tố cơ bản của mô hình đồ thị, một số lý thuyết liên quan
và các thuật toán kinh điển. Còn làm thế nào để xây dựng mô hình đồ thị, vận
dụng các lý thuyết và thuật toán ấy ra sao, nghiên cứu bài toán đồ thị bằng
cách nào, phần lớn đều phải tự hiểu, tự lĩnh hội, rồi kiểm chứng qua thực
hành và nâng cao năng lực bằng quá trình mò mẫm, tổng kết.

Trong quá trình xây dựng mô hình đồ thị, ta thường gặp khó khăn: khó xác định
quan hệ giữa đỉnh, cạnh và trọng số; một số yếu tố quan trọng trong nguyên
mẫu không đưa được vào mô hình hiện có; hoặc mô hình hiện có biểu diễn được
nguyên mẫu nhưng không biết giải ra sao. Để khắc phục những khó khăn ấy, ta
cần dùng một số tư tưởng, phương pháp và kỹ xảo riêng. Bài viết này ghi lại
một vài nhận thức của tôi rút ra trong học tập và thực hành.

\section{Xây dựng mô hình đồ thị}

Trước khi xây dựng mô hình, trước hết ta cần khảo sát toàn diện đối tượng
nghiên cứu, lý tưởng hóa và đơn giản hóa nguyên mẫu; với đề thi, phần lớn
bước này đã được người ra đề hoàn thành. Sau đó ta phân tích sơ bộ nguyên
mẫu, phân biệt các yếu tố và mục tiêu cần giải, rồi sắp xếp quan hệ giữa
chúng. Bước tiếp theo là dùng một mô hình thích hợp để mô tả các yếu tố và
quan hệ ấy.

\subsection{Chọn bỏ yếu tố}

Khi dùng mô hình đồ thị để mô tả đối tượng nghiên cứu, để làm nổi bật các yếu
tố gắn chặt với mục tiêu giải và giảm độ phức tạp khi suy nghĩ, ta không
tránh khỏi phải bỏ bớt một phần yếu tố. Sau đây ta phân tích điều đó qua ví
dụ 1.

\paragraph{Ví dụ 1: Bố trí dây dẫn \textit{Line}.}
Trong mô tả đề, phần trọng điểm liên quan đến các yếu tố như vòng tròn, \(N\)
dây dẫn, \(2N\) đầu mút, quy tắc đánh số, giao cắt giữa các dây. Mục tiêu
cần giải là xây dựng một phương án bố trí dây phù hợp với tình hình giao cắt
đã cho.

Ban đầu, ta chưa quen với cách bố trí dây trong đề. Có lẽ ta có thể vẽ vài
ví dụ vô nghiệm hoặc nhiều nghiệm, nhưng trong kỳ thi không thể dành quá
nhiều thời gian để làm quen với đề. Khi đó, cách duy nhất là nhanh chóng biến
nguyên mẫu xa lạ, khó suy nghĩ thành mô hình quen thuộc và dễ suy nghĩ.

Trước hết phân tích mục tiêu giải: xây dựng một phương án bố trí dây nghĩa là
tìm ra số hiệu hai đầu mút của mỗi dây. Điều kiện mà các số hiệu phải thỏa
chính là tình hình giao cắt giữa các dây.

Tiếp theo, ta phân tích quan hệ giữa số hiệu và giao cắt. Gọi số hiệu hai đầu
mút của dây thứ \(i\) là \(A_i\) và \(B_i\), với \(A_i<B_i\). Quan hệ giữa
hai dây có ba khả năng, tương ứng với ba quan hệ thứ tự sau:
\[
\begin{array}{ll}
  \text{(a) hai dây giao nhau:} & A_1>A_2>B_1>B_2,\\
  \text{(b) không giao nhau:} & A_1>B_1>A_2>B_2,\\
  \text{(c) không giao nhau:} & A_1>A_2>B_2>B_1.
\end{array}
\]
Điểm chung của ba trường hợp là \(A_1>B_1\), \(A_2>B_2\), \(A_1>A_2\), do
quy tắc đánh số quyết định. Điểm khác là: (a) thỏa \(A_2>B_1\) và \(B_1>B_2\);
(b) thỏa \(B_1>A_2\); còn (c) thỏa \(B_2>B_1\). Rõ ràng đây là một quan hệ
thứ tự bộ phận; nói chính xác hơn, nó thỏa phản đối xứng và bắc cầu, nhưng
không tự phản xạ. Nhiệm vụ của ta là từ quan hệ thứ tự bộ phận ấy tìm một
quan hệ thứ tự toàn phần, tức sắp xếp tô pô.

Ta dùng cạnh có hướng trong đồ thị để biểu diễn quan hệ thứ tự bộ phận. Nếu
các cạnh có hướng tạo thành chu trình, bài toán vô nghiệm. Ba trường hợp trên
có thể viết thành các ràng buộc:
\[
\begin{array}{ll}
  \text{(a)} & A_1\to A_2,\quad A_2\to B_1,\quad B_1\to B_2,\\
  \text{(b)} & A_1\to B_1,\quad B_1\to A_2,\quad A_2\to B_2,\\
  \text{(c)} & A_1\to A_2,\quad A_2\to B_2,\quad B_2\to B_1.
\end{array}
\]
Nếu hai dây giao nhau thì dùng (a); nếu không giao nhau thì phải chọn một
trong hai cách (b), (c), tùy cách nào không dẫn đến vô nghiệm.

Nếu chọn (b) mà vô nghiệm, điều đó nghĩa là ngoài cạnh \(B_1\to A_2\), trong
đồ thị đã tồn tại một đường đi từ \(A_2\) đến \(B_1\). Có hai khả năng:
\begin{enumerate}
  \item Có một dây có số hiệu lớn hơn \(2\), chẳng hạn dây \(3\), giao với dây
  \(1\).
  \item Tuy các dây ấy đều không giao với dây \(1\), nhưng trong đó có cặp
  dây không giao đã được biểu diễn bằng dạng (c).
\end{enumerate}
Nếu khả năng thứ nhất xuất hiện, chọn (b) chắc chắn vô nghiệm. Nhưng nếu khả
năng thứ nhất không xuất hiện, trường hợp thứ hai có thể đổi cách biểu diễn:
dùng (b) thay vì (c) cho cặp dây không giao tương ứng, và khi đó sẽ có nghiệm.
Nói cách khác, nếu dây \(i\) và dây \(j\), \(i<j\), không giao nhau, thì việc
có tồn tại dây có số hiệu lớn hơn \(j\) giao với dây \(i\) hay không quyết
định chọn (b) có vô nghiệm hay không.

Tương tự, ta có thể phân tích rằng: nếu dây \(i\) và dây \(j\), \(i<j\),
không giao nhau, thì việc có tồn tại một dãy \(L_1,\ldots,L_m\), trong đó
\(L_1=i\), \(L_m=j\), và dây \(L_{k-1}\) giao với dây \(L_k\) với
\(1<k\le m\), quyết định chọn (c) có vô nghiệm hay không.

Ví dụ, khi dây \(1\) và dây \(3\) không giao, còn dây \(1\) giao dây \(2\) và
dây \(2\) giao dây \(3\), thì \(m=3\), \(L=(1,2,3)\), và cách (c) sẽ vô
nghiệm. Nếu thêm trường hợp dây \(1\) giao dây \(4\), thì với dây \(1\) và
dây \(3\), dùng (b) hay (c) đều vô nghiệm, nên toàn bộ bài toán vô nghiệm.

Phân tích tới đây, thực ra mô hình đồ thị và thuật toán đã hình thành; chi
tiết cài đặt không bàn thêm. Phân tích ở trên nhìn có vẻ dài, nhưng trong
thực tế quá trình suy nghĩ rất tự nhiên. Nguyên nhân là ta đã bỏ đi các yếu
tố trong nguyên mẫu không liên quan tới mục tiêu giải, và chỉ dùng một đồ thị
có hướng để biểu diễn quan hệ giao cắt. So với việc bám chặt vào các chi
tiết như vòng tròn và dây dẫn rồi phân tích trực tiếp nguyên mẫu, độ phức tạp
tư duy khi phân tích mô hình nhỏ hơn nhiều.

Trong quá trình mô hình hóa, bỏ qua các yếu tố ít liên quan đến mục tiêu giải
có thể đơn giản hóa bài toán ở mức thích hợp. Nhưng mọi sự đều có hai mặt:
nếu đơn giản hóa quá mức, mô hình có thể trở nên thiếu chính xác, thậm chí
sai. Vì vậy, chỉ khi tìm được điểm cân bằng giữa ``ngắn gọn'' và ``chính
xác'', ta mới xây dựng được mô hình có hiệu quả tốt nhất.

\subsection{Chọn hệ lý thuyết phù hợp}

Đồ thị gồm ba phần: đỉnh, cạnh và trọng số. Tùy tính chất của ba phần ấy mà
ta có các mô hình đồ thị khác nhau, các lý thuyết và thuật toán khác nhau,
từ đó tạo thành những hệ lý thuyết khác nhau.

Chẳng hạn, đồ thị hai phía chia toàn bộ tập đỉnh \(V\) thành hai tập con và
quy định rằng trong cùng một tập con không có cạnh. Vì vậy, đồ thị hai phía
có các tính chất đặc thù khác đồ thị tổng quát, và thuật toán ghép cặp của nó
cũng đơn giản hơn thuật toán trên đồ thị tổng quát. Ngoài ra còn có cây, đồ
thị có hướng không chu trình, v.v.; chúng thuộc các hệ lý thuyết khác nhau,
có tính chất riêng và thích hợp với các thuật toán giải khác nhau.

Việc đưa trọng số vào lại làm cho mô hình đồ thị và mục tiêu giải đa dạng hơn
nữa. Có trọng số biểu diễn độ dài, thời gian, v.v.; đặc trưng phép toán là
``nối tiếp thì cộng, song song thì lấy trị tối ưu'', tức trọng số của một
đường đi bằng tổng trọng số các cạnh trên đường đi, và mục tiêu thường là tìm
trị tối ưu của trọng số trong toàn đồ thị hoặc giữa hai đỉnh. Có trọng số
biểu diễn dung lượng hoặc lưu lượng; đặc trưng phép toán là ``nối tiếp thì
lấy trị tối ưu, song song thì cộng'', tức trọng số lớn nhất hoặc nhỏ nhất
trên một đường đi quyết định trọng số của cả đường đi, còn mục tiêu là tính
tổng trọng số của các đường đi trong đồ thị hoặc giữa hai đỉnh.

Trên đây chỉ là vài loại trọng số cơ bản. Chúng còn có thể sinh ra các biến
thể, như phép toán trên trọng số mở rộng từ cộng và lấy trị tối ưu sang nhân,
hoặc sang những hàm phức tạp hơn. Có đồ thị không chỉ chứa trọng số cạnh, tức
ánh xạ từ tập cạnh \(E\) đến tập số thực \(R\), mà còn chứa trọng số đỉnh,
tức ánh xạ từ tập đỉnh \(V\) đến \(R\); cũng có đồ thị chứa nhiều loại trọng
số với tính chất khác nhau.

Những khác biệt này tạo nên sự đa dạng của mô hình đồ thị, giúp mô hình đồ
thị thích ứng rộng rãi với nhiều loại bài toán. Nhưng lựa chọn phong phú cũng
làm tăng độ khó của mô hình hóa đồ thị. Với một số đề, ta có thể tự nhiên
liên tưởng đến một mô hình đồ thị nào đó, chẳng hạn thấy biểu thức thì nghĩ
đến cây biểu thức. Với các đề khác, góc phân tích khác nhau sẽ dẫn đến mô
hình khác nhau và hiệu quả khác nhau.

\paragraph{Ví dụ 2: Bò diễu hành \textit{Cows on Parade}.}
Cho một dãy nhị phân độ dài \(2^n+n-1\), yêu cầu dãy chứa mọi xâu nhị phân độ
dài \(n\), tổng cộng \(2^n\) xâu. Khi \(n=3\), một dãy thỏa mãn là
\[
  0001110100,
\]
trong đó các xâu độ dài \(3\) lần lượt là
\[
  000,\ 001,\ 011,\ 111,\ 110,\ 101,\ 010,\ 100.
\]
Bài ``thiết kế trống quay'' tương tự, chỉ khác là yêu cầu sắp thành vòng
tròn chứ không phải thành một đường thẳng. Rõ ràng nếu giải được bài theo
vòng tròn thì cũng giải được bài theo đường thẳng, nên ở đây ta lấy mục tiêu
giải là sắp thành vòng tròn.

\paragraph{Mô hình A.}
Nếu lấy mỗi xâu nhị phân độ dài \(n\) làm một đỉnh, và lấy quan hệ chồng nhau
giữa chúng làm cạnh, ta thu được một đồ thị có hướng. Mục tiêu giải là tìm
một chu trình đi qua mỗi đỉnh đúng một lần: bài toán chu trình Hamilton.

\paragraph{Mô hình B.}
Nếu lấy phần chồng nhau giữa các xâu nhị phân làm đỉnh, và lấy các xâu nhị
phân làm cạnh, ta thu được một đồ thị có hướng khác. Khi \(n=3\), các đỉnh là
\(00,01,10,11\), còn mỗi cạnh là một xâu độ dài \(3\), nối tiền tố độ dài
\(2\) với hậu tố độ dài \(2\). Mục tiêu là tìm một chu trình đi qua mỗi cạnh
đúng một lần: bài toán chu trình Euler.

Trong mô hình B, không khó thấy mọi đỉnh đều liên thông, bậc vào và bậc ra
đều bằng \(2\), nên tồn tại chu trình Euler và cũng không khó tìm. Ngược lại,
bài toán chu trình Hamilton là NP-đầy đủ; nói cách khác, mô hình A không đem
lại thuận lợi nào cho việc giải bài.

Từ đó có thể thấy các mô hình đồ thị thuộc hệ lý thuyết khác nhau rất có thể
tạo ra hiệu quả hoàn toàn khác nhau. Trong ví dụ 2, các mô hình khác nhau
được xây dựng do góc phân tích khác nhau; còn trong ví dụ 3 dưới đây, sự khác
biệt đến từ mức độ phân tích nguyên mẫu sâu hay nông.

\paragraph{Ví dụ 3: Phủ đường đi ít nhất.}
Cho đồ thị có hướng không chu trình \(G(V,E)\). Một phủ đường đi là một tập
đường đi \(P\), sao cho mỗi đỉnh của đồ thị thuộc đúng một đường đi trong
tập. Hãy tìm một phủ đường đi có số đường đi ít nhất.

Với đồ thị mẫu gồm các cạnh \(1\to3\), \(2\to3\), \(3\to4\), \(4\to5\) và
\(2\to5\), ta có \(p_{\min}=2\), với ba phương án:
\[
\{1\to3\to4\to5,\ 2\},\qquad
\{2\to3\to4\to5,\ 1\},\qquad
\{1\to3\to4,\ 2\to5\}.
\]

\paragraph{Cách giải A: mô hình luồng mạng.}
Xây dựng đồ thị mới \(G'\) như sau:
\begin{enumerate}
  \item Tách mỗi đỉnh \(i\in V\) thành hai đỉnh \(i'\) và \(i''\). Thêm cung
  \((i',i'')\), với cận trên dung lượng \(C_{i',i''}=1\) và cận dưới
  \(B_{i',i''}=1\).
  \item Thêm nguồn \(s\) và đích \(t\). Thêm các cung \((s,i')\) và
  \((i'',t)\), với \(C_{s,i'}=C_{i'',t}=+\infty\) và
  \(B_{s,i'}=B_{i'',t}=0\).
  \item Mỗi cung \((i,j)\in E\) được đổi thành cung \((i'',j')\), với
  \(C_{i'',j'}=+\infty\) và \(B_{i'',j'}=0\).
\end{enumerate}
Sau khi xây dựng \(G'\), một đường đi trong \(G\) tương ứng với một đường
luồng trong \(G'\). Vì ta quy định \(C_{i',i''}=B_{i',i''}=1\), đỉnh \(i\)
của đồ thị gốc thuộc đúng một đường đi. Do đó một phủ đường đi ít nhất trong
\(G\) tương ứng với một phương án luồng nhỏ nhất trong \(G'\).

\paragraph{Cách giải B: mô hình đồ thị hai phía.}
Giả sử phủ đường đi là
\[
  P=\{P_1,P_2,\ldots,P_p\},
\]
trong đó \(P_i\), \(1\le i\le p\), là một đường đi trong tập \(P\). Ký hiệu
\(V_i=P_i\cap V\), \(E_i=P_i\cap E\), tức \(V_i\) là tập đỉnh trên \(P_i\) và
\(E_i\) là tập cạnh trên \(P_i\). Rõ ràng \(|V_i|=|E_i|+1\).

Ví dụ trong phương án thứ nhất của hình mẫu,
\[
  P=\{1\to3\to4\to5,\ 2\},
\]
ta có \(V_1=\{1,3,4,5\}\) và
\[
  E_1=\{(1,3),(3,4),(4,5)\}.
\]

Vì \(P\) là một phủ đường đi, nên
\[
  V_1\cup V_2\cup\cdots\cup V_p=V,
  \qquad
  V_i\cap V_j=\varnothing\quad(1\le i,j\le p,\ i\ne j).
\]
Từ đó suy ra
\[
\begin{aligned}
  |V|
  &=\sum_{i=1}^p |V_i|
   =\sum_{i=1}^p (|E_i|+1)
   =p+\sum_{i=1}^p |E_i|,\\
  p&=|V|-\sum_{i=1}^p |E_i|.
\end{aligned}
\]
Lại do \(E_i\cap E_j=\varnothing\) khi \(i\ne j\), đặt
\[
  M=\bigcup_{i=1}^p E_i.
\]
Tập \(M\) cần thỏa điều kiện: mỗi đỉnh trong \(V\) nhiều nhất chỉ liên thuộc
với hai cung trong \(M\), gồm một cung vào và một cung ra. Hơn nữa, chỉ cần
\(M\) thỏa điều kiện này thì nó tương ứng một-một với một phủ đường đi.

Từ phân tích trên, ta xây dựng mô hình đồ thị hai phía:
\begin{enumerate}
  \item Tách mỗi đỉnh \(i\in V\) thành hai đỉnh \(i'\) và \(i''\).
  \item Mỗi cung \((i,j)\in E\) được đổi thành cạnh vô hướng \((i'',j')\).
\end{enumerate}
Trong đồ thị mới \(G'\), đỉnh \(i'\) là đỉnh vào: mọi cung trong \(G\) có đầu
cung là \(i\) đều được đổi thành cạnh nối với \(i'\). Tương tự, đỉnh \(i''\)
là đỉnh ra: mọi cung trong \(G\) có đuôi cung là \(i\) đều được đổi thành
cạnh nối với \(i''\). Khi đó, tập \(M\) trên \(G\) tương ứng với một ghép cặp
trên đồ thị hai phía \(G'\).

Ở trên ta đã phân tích rằng \(M\) và \(P\) tương ứng một-một, và
\(|P|=|V|-|M|\). Vì \(|V|\) là hằng số, chỉ cần tìm ghép cặp cực đại trên đồ
thị hai phía \(G'\) là tìm được phủ đường đi ít nhất của đồ thị gốc \(G\).

So sánh hiệu quả thời gian, không gian và độ phức tạp lập trình của hai cách
giải trên, không khó thấy cách giải B rõ ràng tốt hơn cách giải A. Nguyên
nhân là cách giải B phân tích nguyên mẫu sâu hơn, chỉ trích ra các yếu tố
then chốt để xây dựng mô hình, nên việc giải đơn giản và thuận tiện hơn.

Một điểm đáng chú ý khác là cách giải A dùng thuật toán luồng mạng. Do loại
thuật toán này thường phức tạp khi lập trình, dễ sai, và hệ số thời gian lớn,
biểu hiện của nó trong thi đấu thường không thật nổi bật. Nhưng mô hình luồng
mạng chứa được rất nhiều yếu tố, đặc biệt là nhiều loại trọng số: không chỉ
có trọng số biểu diễn dung lượng và lưu lượng, mà còn có trọng số biểu diễn
chi phí; dung lượng không chỉ có cận trên mà còn có thể có cận dưới. Với
những yếu tố linh hoạt ấy, cộng thêm việc thuật toán luồng mạng nói chung là
thuật toán hiệu quả, mô hình luồng mạng được ứng dụng rất rộng rãi trong giải
quyết vấn đề thực tế và nghiên cứu bài toán đồ thị. Trong các ví dụ sau,
không ít bài cũng dùng mô hình luồng mạng.

Từ các ví dụ trên, ta đã thấy tầm quan trọng của việc chọn hệ lý thuyết phù
hợp. Đôi khi có nhiều mô hình nhìn qua đều đúng để lựa chọn; chỉ khi phân
biệt kỹ những khác biệt nhỏ nhưng then chốt giữa chúng, ta mới chọn được mô
hình thích hợp để giải. Điều này đòi hỏi ta vừa quen thuộc với thuật toán và
lý thuyết đồ thị, vừa hiểu rõ nguyên mẫu.

Dù là quyết định chọn bỏ yếu tố hay đứng trước lựa chọn giữa các hệ lý thuyết
khác nhau, nguyên tắc chung đều là: chính xác, rõ ràng, ngắn gọn.

\begin{itemize}
  \item \term{Chính xác} nghĩa là mô hình phải trung thành với nguyên mẫu,
  không bóp méo tính chất vốn có của nguyên mẫu, cũng không ngầm giả định
  nguyên mẫu có những tính chất chưa được chứng minh.
  \item \term{Rõ ràng} nghĩa là giống như chương trình có tính dễ đọc, mô
  hình đồ thị cũng có độ rõ ràng. Lấy mô hình có cấu trúc rõ làm đối tượng
  phân tích sẽ giảm mạnh độ phức tạp khi suy nghĩ.
  \item \term{Ngắn gọn} nghĩa là các yếu tố trong mô hình đơn giản, sáng sủa,
  không rối rắm, và mô hình thuận tiện cho phân tích, giải. Yêu cầu này xuất
  phát từ việc nhiều thuật toán và lý thuyết trong đồ thị khá sâu và phức
  tạp, khó vận dụng và biến đổi, nên cố gắng tránh dùng khi không cần.
\end{itemize}

\section{Chuyển hóa mô hình đồ thị}

Khi xây dựng mô hình đồ thị để giải bài, ta thường gặp tình huống: sau khi
dễ dàng xây dựng một mô hình, ta lại không biết phân tích hay giải như thế
nào, hoặc phát hiện mô hình đã xây dựng không thể biểu diễn một tính chất rất
quan trọng của nguyên mẫu. Khi đó ta có buộc phải bỏ mô hình hiện có không?
Không. Nếu vận dụng một số phương pháp và kỹ xảo, mô hình hiện có rất có thể
được chuyển hóa thành một mô hình mới vừa mô tả chính xác nguyên mẫu vừa dễ
giải.

Trong các phương pháp và kỹ xảo ấy, có những cách chỉ thích hợp với một mô
hình đồ thị cụ thể, chẳng hạn kỹ xảo chuyển cây nhiều nhánh thành cây nhị
phân để xử lý. Ngoài ra còn có các phương pháp và kỹ xảo có phạm vi tác dụng
rộng hơn. Phần dưới sẽ bàn về một trong số đó: chuyển hóa bằng tách.

Chuyển hóa bằng tách có thể chia thành ba loại: đỉnh thành cạnh, đỉnh thành
đỉnh, và cạnh thành cạnh. Loại thứ nhất rất thường dùng nhưng cách tách khá
đơn giản; loại thứ hai không chỉ thường dùng mà còn linh hoạt, đa dạng; cách
tách và tác dụng của loại thứ ba tương tự loại thứ hai. Vì vậy, phần dưới
trình bày kỹ loại thứ hai, còn loại thứ nhất và thứ ba chỉ nói ngắn gọn.

\subsection{Đỉnh thành cạnh}

``Tách đỉnh thành cạnh'' là một loại rất thường dùng trong chuyển hóa bằng
tách. Khi cần làm cho đỉnh trong mô hình hiện có mang tính chất của cạnh,
tách đỉnh thành cạnh là một cách tốt.

Trong ví dụ 3, khi xây dựng mô hình A, ta tách mỗi đỉnh thành một cạnh có
dung lượng với cận trên và cận dưới để thỏa yêu cầu ``mỗi đỉnh được đi qua
đúng một lần''. Ngoài ra, khi giải luồng mạng có dung lượng trên đỉnh, để thu
được mô hình luồng mạng chuẩn, ta thường cũng tách đỉnh thành cạnh, chuyển
ràng buộc dung lượng trên đỉnh sang cạnh tương ứng.

Trong ứng dụng thực tế, có nhiều ví dụ tách đỉnh thành cạnh, và cách tách
phần lớn tương tự nhau: tách mỗi đỉnh \(i\) thành hai đỉnh \(i'\) và \(i''\);
mọi cung có đầu cung là \(i\) được đổi thành có đầu cung là \(i'\); mọi cung
có đuôi cung là \(i\) được đổi thành có đuôi cung là \(i''\). Trọng số và
thuộc tính cũ trên các cung ấy được giữ hay đổi tùy tình huống. Sau đó thêm
một cung từ \(i'\) đến \(i''\), và gán cho cung này thuộc tính vốn thuộc về
đỉnh \(i\).

\subsection{Đỉnh thành đỉnh}

Đây là loại thứ hai của chuyển hóa bằng tách: tách một đỉnh thành nhiều đỉnh.
Nó chủ yếu dùng để tách nhiều tính chất khác nhau của cùng một đỉnh ra, gán
riêng cho các đỉnh mới, để mỗi đỉnh có tính chất đơn nhất, rồi thông qua một
số cải tạo để xử lý thống nhất chúng.

Chẳng hạn trong ví dụ 3, khi xây dựng mô hình B, để biểu diễn rằng một đỉnh
bất kỳ \(i\) nhiều nhất được liên thuộc với một cung vào và một cung ra, ta
tách \(i\) thành một đỉnh vào \(i'\) và một đỉnh ra \(i''\). Bằng cách nối
mọi cung vào với \(i'\), và mọi cung ra với \(i''\), mỗi đỉnh mới chỉ cần
liên thuộc với một cạnh, nên có thể dùng thuật toán ghép cặp để xử lý thống
nhất.

Cách tách đỉnh rất đa dạng, không chỉ có dạng trên. Ví dụ sau dùng một cách
tách khác với ví dụ 3.

\paragraph{Ví dụ 4: Máy dò vật cản \textit{Mars explorer}.}
Trong hình chữ nhật \(P\times Q\), địa hình của mỗi ô có thể thuộc một trong
ba loại:
\begin{enumerate}
  \item Vật cản: không thể đi qua.
  \item Đất bằng: bằng phẳng, không có vật, có thể đi qua.
  \item Đá: có thể đi qua; lần đầu đi qua thì nhặt được một viên đá.
\end{enumerate}
Cần tìm \(M\) đường đi, với \(M\) cho trước, sao cho tổng số ô đá được các
đường đi phủ tới là lớn nhất. Các đường đi phải thỏa:
\begin{enumerate}
  \item Mỗi đường đi đều từ \((1,1)\) đến \((P,Q)\), mỗi bước chỉ được đi
  sang đông hoặc xuống nam một ô.
  \item Đường đi không đi qua vật cản.
\end{enumerate}

\paragraph{Xây dựng mô hình luồng mạng.}
Theo suy nghĩ thông thường, xem mỗi đường đi là một đường luồng có lưu lượng
\(1\), rồi dựa trên hai điều kiện của đường đi để xây dựng mô hình:
\begin{enumerate}
  \item Xem mỗi ô là một đỉnh, và dẫn hai cung lần lượt đến ô kề phía đông và
  phía nam nếu chúng tồn tại. Lưu lượng trên cung biểu diễn số đường đi đi qua
  cung ấy, nên dung lượng là \(+\infty\), hoặc chỉ cần không nhỏ hơn \(M\).
  \item Xóa các cung liên quan đến đỉnh vật cản.
\end{enumerate}
Đây chỉ là mô hình sơ bộ, vì rõ ràng nó chưa phân biệt ``đất bằng'' và
``đá''. Khó khăn chính là: đá chỉ có thể nhặt lần đầu khi đi qua, sau đó ô
này trở thành đất bằng. Vì thế cần tách đỉnh \(i\) có đá thành hai đỉnh:
\(i'\) và \(i''\), lần lượt biểu diễn trạng thái ``đá'' và ``đất bằng''. Cách
cải tạo cụ thể như sau:
\begin{enumerate}
  \item \(i''\) kế thừa mọi cung vào và cung ra liên quan đến \(i\), dung
  lượng không đổi, lợi ích bằng \(0\).
  \item Dẫn một cung từ \(i''\) đến \(i'\), đặt dung lượng bằng \(1\), lợi ích
  bằng \(1\).
  \item \(i'\) kế thừa mọi cung ra liên quan đến \(i\), dung lượng đổi thành
  \(1\), hoặc lớn hơn \(1\) cũng được, lợi ích bằng \(0\).
\end{enumerate}
Sau chuyển hóa trên, mục tiêu của bài toán gốc tương ứng với luồng có lợi ích
lớn nhất và lưu lượng bằng \(M\) trên mô hình mới.

Hai ví dụ trên đều tách một đỉnh thành hai đỉnh, nhưng điều này hoàn toàn
khác tách đỉnh thành cạnh. Ở đây, mục đích là tách các tính chất khác nhau
của một đỉnh ra; cần tách bao nhiêu loại tính chất thì tách thành bấy nhiêu
đỉnh. Còn tách đỉnh thành cạnh là để đỉnh mang tính chất của cạnh, nên luôn
cố định tách một đỉnh thành hai đỉnh và thêm cạnh giữa hai đỉnh mới.

\subsection{Cạnh thành cạnh}

Đây là loại thứ ba của chuyển hóa bằng tách: tách một cạnh thành nhiều cạnh.
Nó giống loại thứ hai ở chỗ cũng dùng để tách yếu tố. Dưới đây chỉ nêu một
ví dụ ngắn.

\paragraph{Ví dụ 5: Luồng mạng có dung lượng với cận trên và cận dưới.}
Đây là một bài toán luồng mạng chuẩn. Nhiều sách đã có cách giải và chứng
minh liên quan, nên ở đây chỉ trình bày ngắn quá trình chuyển đồ thị gốc
thành mô hình luồng cực đại thông thường:
\begin{enumerate}
  \item Thêm hai đỉnh mới: nguồn phụ \(s'\) và đích phụ \(t'\), làm nguồn và
  đích mới.
  \item Tách cung \((i,j)\) của đồ thị gốc thành ba cung mới:
  \((i,j)\), \((s',j)\), \((i,t')\). Đặt
  \[
    C'_{i,j}=C_{i,j}-B_{i,j},\qquad
    C'_{s',j}=C'_{i,t'}=B_{i,j},
  \]
  trong đó \(C_{i,j}\), \(B_{i,j}\) lần lượt là cận trên và cận dưới dung
  lượng của cung \((i,j)\) trong đồ thị gốc, còn \(C'_{i,j}\),
  \(C'_{s',j}\), \(C'_{i,t'}\) là cận trên dung lượng của các cung mới.
  \item Thêm hai cung mới dung lượng \(+\infty\) giữa nguồn và đích cũ:
  \((s,t)\) và \((t,s)\).
\end{enumerate}
Trong ba bước trên, bước thứ hai là quan trọng nhất: nó tách cận trên và cận
dưới của dung lượng trên đồ thị gốc.

Đi xa hơn, nếu tính chất của đỉnh hoặc cạnh thay đổi theo thời gian, ta cũng
có thể dùng cách tương tự loại thứ hai hoặc thứ ba để tách chúng thành nhiều
đỉnh hoặc cạnh, biểu diễn tính chất ở các thời điểm khác nhau. Ví dụ 6 cho
thấy dạng chuyển hóa này được áp dụng như thế nào.

\paragraph{Ví dụ 6: Gia viên \textit{Homeland}.}
Sau khi phân tích kỹ đề, không khó nghĩ đến việc xây dựng mô hình luồng mạng:
\begin{enumerate}
  \item Lấy trạm vũ trụ làm đỉnh.
  \item Lấy tàu vũ trụ qua lại giữa các trạm làm cạnh.
  \item Lấy sức chứa tối đa của tàu làm cận trên dung lượng.
\end{enumerate}
Mục tiêu là trong thời gian ngắn nhất đưa một lượng luồng cố định, tức con
người, từ nguồn là Trái Đất đến đích là Mặt Trăng. Nhưng ở đây, thời gian
không giống chi phí; hơn nữa cạnh tồn tại giữa những cặp đỉnh khác nhau theo
thời điểm \(t\). Do đó tách đỉnh theo thời gian là cần thiết: tách mỗi đỉnh
thành \(t\) đỉnh, rồi phân phối cạnh vào các cặp đỉnh mới tương ứng theo thời
điểm khác nhau.

Như vậy, mô hình vốn biến đổi động được chuyển thành một mô hình mới tĩnh.

Từ các ví dụ trên có thể thấy: tính chất cần tách khác nhau thì cách tách
đỉnh hoặc cạnh cũng khác nhau. Chỉ khi phân tích chính xác tính chất mà đỉnh
hoặc cạnh cần biểu diễn, ta mới chuyển hóa được thành mô hình phù hợp.

Qua phân tích và tổng kết ba loại chuyển hóa bằng tách đối với đỉnh và cạnh,
không khó nhận ra: mục đích và phạm vi ứng dụng của chúng khác nhau, nhưng
phương pháp đều nằm ở chữ ``tách''. Khi áp dụng, ta hoàn toàn không cần câu
nệ hình thức cụ thể; miễn là mô hình hóa và giải bài cần, ta có thể tách theo
nhu cầu. Chỉ cần tách đúng, ba loại chuyển hóa trên có thể được dùng rộng hơn
nhiều.

\subsection{Chuyển hóa bằng phần bù}

Ngoài chuyển hóa bằng tách, chuyển hóa bằng phần bù cũng là một phương pháp
chuyển hóa không phụ thuộc vào mô hình cụ thể.

Trong phân tích ví dụ 4, ta còn để lại một vấn đề: chuyển luồng có lợi ích
lớn nhất với lưu lượng cố định thành luồng chi phí nhỏ nhất để giải. Việc này
đòi hỏi sửa trọng số trên cạnh: đổi chi phí của mọi cung giữa các cặp đỉnh
biểu diễn ``đất bằng'' thành \(1\); đổi chi phí của mọi cung chỉ vào đỉnh
``đá'' hoặc đi ra từ những đỉnh này thành \(0\). Khi đó, lợi ích của mỗi
đường luồng trong mô hình ban đầu bằng \(P+Q-2\), tức độ dài đường đi, trừ đi chi
phí trên đường đi. Như vậy ta đã chuyển mục tiêu lợi ích lớn nhất thành chi
phí nhỏ nhất.

Kỹ xảo chuyển hóa bằng cách lấy một giá trị cố định trừ đi trọng số cũ như
vậy được gọi là chuyển hóa phần bù của trọng số.

Đã có chuyển hóa phần bù của trọng số thì tự nhiên cũng có chuyển hóa phần bù
của đỉnh và cạnh. Chuyển hóa phần bù của đỉnh và cạnh là kỹ xảo lấy tập toàn
phần, như tập đỉnh, tập cạnh của đồ thị hoặc tập khác, trừ đi một tập cụ thể,
như tập đề bài cho hoặc tập mục tiêu cần tìm. Dưới đây ta phân tích bằng ví
dụ cụ thể.

\paragraph{Ví dụ 7: Bài toán tập độc lập lớn nhất và các bài toán tương đương.}
Bài toán tập độc lập lớn nhất được mô tả như sau: cho đồ thị vô hướng
\(G(V,E)\), trong đó \(A\subseteq V\), và
\[
  E\cap \{(i,j)\mid i,j\in A\}=\varnothing.
\]
Hãy tìm \(A\) sao cho \(|A|\) lớn nhất. Bài toán này có hai bài toán tương
đương:

\paragraph{1. Bài toán phủ nhỏ nhất.}
Thông thường, ta chuyển tập độc lập lớn nhất thành tập phủ nhỏ nhất, rồi giải
dựa trên các luật logic. Đặt \(B=V-A\). Khi đó \(B\subseteq V\), và với mọi
cạnh \((i,j)\in E\), có \(i\in B\) hoặc \(j\in B\), nên \(B\) là tập phủ nhỏ
nhất của đồ thị \(G(V,E)\). Trong quá trình chuyển hóa này, ta dùng chuyển
hóa phần bù của đỉnh: lấy tập đỉnh \(V\) trừ đi tập mục tiêu \(A\) để thu
được tập mục tiêu mới \(B\).

\paragraph{2. Bài toán clique lớn nhất.}
Mô tả bài toán: cho đồ thị vô hướng \(G(V,E)\), trong đó \(C\subseteq V\).
Với mọi hai đỉnh \(i,j\in C\), \(i\ne j\), đều có \((i,j)\in E\). Hãy tìm
\(C\) sao cho \(|C|\) lớn nhất.

Khi ta đã biết và có thể chứng minh bài toán này là NP-đầy đủ, chỉ cần chuyển
nó thành bài toán tập độc lập lớn nhất để giải, ta cũng chứng minh được bài
toán sau là NP-khó. Bước chuyển hóa này dùng chuyển hóa phần bù của cạnh:
đặt
\[
  U=\{(i,j)\mid i,j\in V\},\qquad E'=U-E.
\]
Trong đồ thị vô hướng \(G'(V,E')\), ta có \(C\subseteq V\) và
\[
  E'\cap\{(i,j)\mid i,j\in C\}=\varnothing,
\]
nên \(C\) cũng là tập độc lập lớn nhất của \(G'\).

Chẳng hạn, trong một cặp đồ thị bù nhau \(G'\) và \(G\), các đỉnh đen trên
\(G'\) tạo thành tập độc lập lớn nhất, còn các đỉnh trắng là tập phủ nhỏ
nhất. Đồ thị \(G\) có cùng tập đỉnh với \(G'\), còn tập cạnh \(E\) là phần bù
của \(E'\); khi đó các đỉnh đen tạo thành tập đỉnh của clique lớn nhất trong
\(G\).

Từ các ví dụ trên có thể thấy chuyển hóa phần bù thường là chuyển hóa tương
đương, và thường làm mục tiêu giải thay đổi ở mức nhất định. Dù là chuyển hóa
phần bù của đỉnh, cạnh hay trọng số, mục đích đều là làm mô hình chuyển theo
hướng thuận tiện hơn cho suy nghĩ và giải.

Ngoài hai phương pháp đã phân tích là chuyển hóa bằng tách và chuyển hóa bằng
phần bù, còn có rất nhiều phương pháp và kỹ xảo chuyển hóa mô hình đồ thị
khác. Chúng cũng có thể được dùng kết hợp với nhau, giống như trong ví dụ 4:
sau vài lần chuyển hóa mô hình, cuối cùng ta thu được mô hình đồ thị cần có.

\section{Kết luận}

Bài viết này chủ yếu phân tích hai điểm then chốt trong xây dựng mô hình đồ
thị và một số phương pháp, kỹ xảo chuyển hóa mô hình đồ thị. Trong quá trình
mô hình hóa thực tế, chúng không thể tách rời nhau: trích xuất đúng các yếu
tố trong nguyên mẫu; tương ứng chúng với các phần tử trong một hệ lý thuyết
cụ thể; xây dựng mô hình sơ bộ; rồi chuyển hóa thích hợp theo nhu cầu. Tới
đây, một mô hình đồ thị thích hợp để giải mới được xây dựng thành công.

Trong toàn bộ quá trình ấy, dù là nắm nguyên tắc mô hình hóa hay vận dụng
phương pháp chuyển hóa mô hình, tất cả đều tuân theo một điểm: tính chất của
nguyên mẫu quyết định mô hình. Nếu gò nguyên mẫu vào một mô hình không thích
hợp, thường ta lại phá hỏng tính chất then chốt của nguyên mẫu. Khi đó, cho
dù mô hình được xây dựng tinh xảo và kinh điển đến đâu, nó cũng khó vượt qua
kiểm nghiệm.

Thuật toán và lý thuyết đồ thị rất độc đáo, tinh diệu, nhưng khó vận dụng
tùy ý; việc xây dựng và chuyển hóa mô hình đồ thị rất linh hoạt, nên cũng khó
nắm vững. Vì vậy, nghiên cứu mô hình đồ thị không phải việc một sớm một
chiều, mà cần bền bỉ lâu dài. Bài viết này chỉ là một vài nhận thức nông cạn
của cá nhân tôi về mô hình hóa đồ thị, và cũng chỉ là một khởi đầu. Tôi tin
rằng khi nhận thức tiếp tục sâu hơn và trí tuệ tập thể được phát huy, mô hình
đồ thị nhất định sẽ có những ứng dụng rộng lớn và đặc sắc hơn.

\appendix

\section{Ghi chú phụ lục và tài liệu tham khảo}

Trọng tâm của bài viết là mô hình hóa đồ thị; việc giải, giải thích và kiểm
nghiệm mô hình không nằm trong phạm vi thảo luận chính. Các thuật toán và lý
thuyết đồ thị được nhắc tới trong bài đều được giới thiệu chi tiết trong các
sách dưới đây, nên trong bài chỉ phân tích ngắn.

\subsection*{Tài liệu tham khảo}

\begin{enumerate}
  \item Wu Wenhu, Wang Jiande, \emph{Phân tích thuật toán thực dụng và thiết
  kế chương trình}, Nhà xuất bản Công nghiệp Điện tử, 1998.
  \item Ren Shanqiang, Lei Ming, \emph{Mô hình toán học}, Nhà xuất bản Đại
  học Trùng Khánh, 1998.
  \item Pan Jingui, Gu Tiecheng và các tác giả khác, \emph{Cấu trúc dữ liệu
  và thuật toán thông dụng hiện đại trong máy tính}, Nhà xuất bản Đại học Nam
  Kinh, 1994.
  \item Wu Wenhu, Wang Jiande, \emph{Phân tích đề thi Olympic Tin học thanh
  thiếu niên quốc tế và trong nước, 1994--1995}, Nhà xuất bản Đại học Thanh
  Hoa, 1997.
  \item Wu Wenhu, Zhao Peng, \emph{Đề thi và phân tích cuộc thi lập trình máy
  tính Hoa Kỳ 1993--1996}, Nhà xuất bản Đại học Thanh Hoa, 1999.
  \item Wu Wenhu, Wang Jiande, \emph{Hướng dẫn Olympic Tin học thanh thiếu
  niên quốc tế và toàn quốc: thuật toán đồ thị và thiết kế chương trình}, Nhà
  xuất bản Đại học Thanh Hoa, 1997.
  \item Tài liệu đội tuyển quốc gia năm 1999.
  \item Xie Zheng, Li Jianping, \emph{Thuật toán mạng và lý thuyết độ phức
  tạp}, Nhà xuất bản Đại học Khoa học Kỹ thuật Quốc phòng, 1995.
  \item Ni Zhaozhong, \emph{Tuyển tập bài viết bài tập của đội tuyển}, 1998.
\end{enumerate}

\begin{flushright}
Xu Jing, An Huy\\
2000-01-15
\end{flushright}

\section{Phụ lục: Bố trí dây dẫn \textit{Line}}

\paragraph{Mô tả bài toán.}
Phòng thí nghiệm ACM gần đây đang thử nghiệm một thiết bị chuyển tiếp dây dẫn
kiểu mới. Thiết bị này là một hộp tròn, trên chu vi có \(2N\) điểm chuyển
tiếp phân bố đều. Hiện có \(N\) dây dẫn cần nối vào thiết bị. Do yêu cầu đặc
biệt, một số dây bắt buộc phải giao nhau. Vì số lượng dây quá nhiều, tình
hình quá phức tạp, con người khó làm thủ công, nên hãy viết chương trình giúp
nhân viên bố trí các dây dẫn.

Trên một vòng tròn có \(2N\) điểm, lần lượt đánh số
\[
  1,2,3,\ldots,2N-1,2N.
\]
Lấy \(2N\) điểm này làm đầu mút có thể nối được \(N\) dây dẫn, mỗi đầu mút
chỉ được nối với một dây. Mỗi dây đều có số hiệu riêng; số hiệu dây được sắp
theo số hiệu nhỏ hơn trong hai đầu mút của dây. Đã biết tình hình giao cắt
giữa các dây. Hãy tìm một cách có thể để bố trí các dây sao cho phù hợp yêu
cầu.

\paragraph{Yêu cầu vào ra.}
Dữ liệu vào nằm trong tệp \texttt{Input.txt}.
\begin{itemize}
  \item Dòng 1 là tổng số dây \(N\), với \(1\le N\le 400\).
  \item Dòng 2 là tổng số giao điểm \(M\), với \(1\le M\le N(N-1)\).
  \item Các dòng \(3\) đến \(M+2\), mỗi dòng ghi số hiệu hai dây giao nhau;
  số hiệu dây đứng trước nhỏ hơn số hiệu dây đứng sau. Các giao điểm này đã
  được sắp xếp: giao điểm có số hiệu dây thứ nhất nhỏ hơn đứng trước; nếu số
  hiệu dây thứ nhất bằng nhau thì giao điểm có số hiệu dây thứ hai nhỏ hơn
  đứng trước.
\end{itemize}

Dữ liệu ra nằm trong tệp \texttt{Output.txt}. Ghi \(N\) dòng, mỗi dòng hai số,
lần lượt là số hiệu hai đầu mút của một dây.

\paragraph{Dữ liệu mẫu.}
\begin{verbatim}
Input.txt
4
4
12
13
24
34

Output.txt
14
27
36
58
\end{verbatim}

\paragraph{Chấm điểm.}
Trong thời gian quy định, nếu chương trình thoát bình thường và thu được đáp
án đúng thì nhận điểm của bộ thử đó. Giới hạn thời gian dài nhất của mọi bộ
thử không vượt quá \(15\) giây.

\section{Phụ lục: Gia viên \textit{Homeland}}

\subsection{Ý nghĩa đề bài và phân tích thuật toán}

Sau khi phân tích kỹ đề, ta không khó nhận ra bài này có đủ các yếu tố của
một mạng luồng:
\begin{itemize}
  \item nguồn: Trái Đất, ký hiệu \(0\);
  \item đích: Mặt Trăng, ký hiệu \(-1\);
  \item đỉnh trung gian: các trạm vũ trụ \(1,\ldots,n\);
  \item cạnh: các tàu vũ trụ đi lại giữa các trạm;
  \item dung lượng cạnh: số người tối đa mà tàu vũ trụ có thể chở.
\end{itemize}

Điểm khác là bài này có thêm yếu tố thời gian, và cho trước lưu lượng, yêu
cầu thời gian ngắn nhất. Khó khăn là: mạng luồng thông thường không thể biểu
diễn dung lượng và lưu lượng khác nhau giữa hai đỉnh ở các thời điểm khác
nhau. Vì vậy cần cải tạo mạng luồng bằng cách tách đỉnh.

Nếu thời gian ngắn nhất là \(T\), ta tách mỗi đỉnh \(V_i\), gồm cả nguồn và
đích, thành \(T+1\) đỉnh
\[
  V_{i,0},V_{i,1},\ldots,V_{i,T}.
\]
Như vậy, hành động tàu vũ trụ rời \(V_i\) ở thời điểm \(t-1\) và đến
\(V_j\) ở thời điểm \(t\), với \(t\le T\), có thể được biểu diễn bằng cạnh
\[
  V_{i,t-1}\to V_{j,t}.
\]
Nếu \(C(V_i,t-1,V_j,t)=a\) và \(f(V_i,t-1,V_j,t)=b\), điều đó nghĩa là tàu có
thể chở \(a\) người và lúc ấy đang chở \(b\) người. Vì mỗi trạm vũ trụ có thể
cho số người không hạn chế ở lại, với mỗi \(V_{i,t-1}\) đều cần có một cạnh
dung lượng không hạn chế trỏ tới \(V_{i,t}\).

Sau cải tạo này, nguồn là đỉnh \(V_{0,0}\), đích là đỉnh \(V_{-1,T}\). Lưu
lượng từ nguồn đến đích biểu diễn số người nhiều nhất có thể đưa từ Trái Đất
đến Mặt Trăng trước thời điểm \(T\). Ta thử các giá trị \(T\) từ nhỏ đến lớn;
thời điểm đầu tiên có thể đưa toàn bộ \(k\) người tới Mặt Trăng chính là lời
giải tối ưu.

Trong cài đặt, ta dùng thuật toán tìm đường tăng bằng tìm kiếm theo chiều
rộng, tức Edmonds--Karp, để tìm luồng cực đại. Hơn nữa, mỗi khi tăng \(T\)
lên \(1\), ta chỉ cần tiếp tục tăng luồng trên cơ sở lưu lượng đã có, không
cần giải lại từ đầu.

Độ phức tạp không gian là \(O(Tn^2)\). Vì \(T\) có thể lớn hơn \(1000\), giá
trị lớn nhất hiện tìm được là \(1159\), trong Turbo Pascal 7 chỉ lưu được
mạng lưu lượng; mạng dung lượng cần được dựng lại trong quá trình giải, may
là việc này không quá phiền.

Độ phức tạp thời gian là \(O((nT)^3)\). Trong toàn bộ quá trình giải, số lần
tìm kiếm rộng không tìm được đường tăng là \(O(T)\), vì hễ không tìm được
đường tăng thì \(T\) tăng thêm \(1\). Số lần tìm được đường tăng là
\(O(|V||E|)\), còn một lần tìm kiếm rộng có độ phức tạp \(O(|E|)\). Tổng hợp
các yếu tố này, thay \(|E|=(m+n)T\), trong đó \(m\) là số tàu vũ trụ và \(n\)
là số cạnh lưu người ở lại từ \(V_{i,t-1}\) đến \(V_{i,t}\), cùng
\(|V|=nT\), ta được \(O(T^3n(m^2+n^2))\). Với phạm vi đề bài
\(n\le20\), \(m\le13\), có thể xem \(m=O(n)\), nên độ phức tạp thời gian viết
gọn là \(O((nT)^3)\).

\subsection{Một phỏng đoán}

Trong quá trình tạo dữ liệu thử, tôi phát hiện lời giải tối ưu \(T\) tăng có
quy luật theo số người \(k\): khi \(k\) tăng thêm một giá trị cố định \(d\),
\(T\) tăng thêm một giá trị cố định \(t\). Viết thành biểu thức:
\[
  T=t\left\lfloor\frac{k+a}{d}\right\rfloor+b,
\]
trong đó \(t>0\), \(d>0\), \(a\in[0,d)\), còn \(b\) không bị hạn chế; các giá
trị \(t,d,a,b\) do cấu trúc và dung lượng của mạng luồng quyết định.

Chẳng hạn với một dữ liệu mẫu trong bản gốc, mạng luồng tương ứng có
\[
  t=4,\qquad d=6,\qquad a=5,\qquad b=-1,
\]
tức
\[
  T=4\left\lfloor\frac{k+5}{6}\right\rfloor-1.
\]
Bảng đáp án đi kèm thể hiện các đoạn:
\[
  k=1\ldots6\Rightarrow T=3,\quad
  k=7\ldots12\Rightarrow T=7,\quad
  k=13\ldots18\Rightarrow T=11,\quad
  k=19\ldots24\Rightarrow T=15,\ldots
\]
Trên đây chỉ là một hiện tượng tôi phát hiện trong thực hành. Do trình độ lý
thuyết còn hạn chế, tôi chưa thể giải thích hoặc chứng minh. Nhưng nếu nó
đúng, ít nhất ta có thể trước hết giải các trường hợp \(k\) nhỏ, tốn rất ít
thời gian, tìm ra \(t,d,a,b\), rồi tính lời giải khi \(k\) rất lớn; thậm chí
có lẽ chỉ dựa vào dữ liệu vào cũng có thể tính trực tiếp.

\subsection{Chương trình và chú thích}

\begingroup
\scriptsize
\begin{verbatim}
{$A+,B-,D+,E-,F-,G+,I+,L+,N+,O-,P-,Q-,R-,S-,T+,V-,X+}
{$M 16384,0,655360}
program Escape;
const
    Infns='Homeland.in';
    Outfns='Homeland.out';
    LimitM=13;
    LimitN=20;
    LimitK=50;
    LimitTime=1300; { Gia tri thoi gian lon nhat }

       LimitNode=500; { Dung cho tim kiem rong trong thuat toan max-flow }
       LImitDanger=LimitNode-50;
type
       TMap=array[0..LimitN,-1..LimitN] of byte;
       TTimeMap=array[0..LimitTime] of ^TMap;
       TStationBool=array[-1..LimitN] of boolean;
       TShip=record
                content,moor:byte;
                mooring:array[1..LimitN+2] of shortint;
             end;
       TSpaceShip=array[1..LimitM] of TShip;
       TNode=record
                delta,station:shortint;
                level:word;
             end;
       TLinearNode=array[1..LimitNode] of TNode;
       TStationInt=array[0..LimitTime,-1..LimitN] of shortint;
var
       Ship      :TSpaceShip;    { Tau vu tru }
       Flow           :TTimeMap;    { Luu luong tren moi canh o moi thoi diem }
       Father :TStationInt; { Ghi nut cha cua moi nut tren duong tang }
       OpenList:^TLinearNode; { Hang doi cho tim kiem rong trong max-flow }
       n,m,k     :byte;
       pre,now,
       head,tail :word;       { Con tro dau va cuoi hang doi }

procedure Init;      { Doc du lieu vao }
 var
     f:Text;
     i,j:byte;
 begin
     assign(f,Infns); reset(f);
     readln(f,n,m,k);
     for i:=1 to m do
       with Ship[i] do
           begin
            read(f,content,moor);
            for j:=1 to moor do { Doc thong tin tau vu tru }
                read(f,mooring[j]);
            readln(f);
           end;
     close(f);
 end;

procedure Out(totaltime:word); { Xuat ket qua }
 var
     f:text;
 begin
     assign(f,Outfns);rewrite(f);
     writeln(f,totaltime);
    close(F);halt;
 end;

function NoAnswer:boolean;           { Kiem tra truong hop vo nghiem }
  var
      i,j:byte;
      flag:boolean;
      Touch:TStationBool;
  begin
      fillchar(Touch,sizeof(Touch),false);
      Touch[0]:=true;
      repeat
        flag:=true;
        for i:=1 to m do
            with Ship[i] do
              for j:=1 to moor do
                  if Touch[mooring[j]] and (not Touch[mooring[j mod moor +1]])
                    then begin
                            Touch[mooring[j mod moor +1]]:=true;
                            flag:=false;
                         end;
      until flag;
      NoAnswer:=not Touch[-1];
  end;

function min(a,b:byte):byte;     { Tra ve gia tri nho hon cua a va b }
  begin
     if a<=b then min:=a
               else min:=b;
  end;

procedure IncFlow(delta:shortint;level:word);
{ Tang luong theo duong tang da tim duoc; delta la luong tang them }
  var
      i,fa:shortint;
  begin
      i:=-1;      dec(k,delta); { Them delta nguoi da duoc sap xep di chuyen }
      repeat
        fa:=Father[level,i];
        if fa>=0
             then begin { Tang luong tren cung thuan }
                       inc(Flow[level-1]^[fa,i],delta);
                       i:=fa;   dec(level);
                   end
             else begin      { Tang luong tren cung nguoc }
                       dec(Flow[level]^[i,-fa],delta);
                       i:=-fa;  inc(level);
                   end;
      until i=0;
  end;

function Open:boolean;
{ Mo rong nut head; Open=true nghia la chua tim thay duong tang,
  can tiep tuc mo rong nut. }
  var
      Map:array[-1..LimitN] of shortint;
      i,which,t:shortint;
      nextlevel:word;
  begin
  with OpenList^[head] do
      begin
        fillchar(Map,sizeof(Map),0);
        Map[station]:=LimitK;
        for i:=1 to m do { Tinh dung luong cac cung xuat phat tu (level,station) }
             with Ship[i] do
               begin
                    which:=level mod moor +1;
                    if mooring[which]=station
                      then inc(Map[mooring[which mod moor +1]],content);
               end;
        if Map[-1]>Flow[level]^[station,-1] then { Co the tang luong }
             begin
               Father[level+1,-1]:=station;
               IncFlow(min(delta,Map[-1]-Flow[level]^[station,-1]),level+1);
               Open:=false;         exit;
             end;
        Open:=true;
        nextlevel:=level+1;
        if level<pre then
             for i:=0 to n do           { Tim cung thuan }
               if (Father[nextlevel,i]=100) and (Map[i]>Flow[level]^[station,i]) then
                    begin
                      Father[nextlevel,i]:=station;
                      inc(tail);
                      OpenList^[tail].delta:=min(delta,Map[i]-Flow[level]^[station,i]);
                      OpenList^[tail].level:=nextlevel;
                      OpenList^[tail].station:=i;
                    end;
        if level=0 then exit;
        nextlevel:=level-1;
        for i:=1 to n do { Tim cung nguoc }
             if (Father[nextlevel,i]=100) and (Flow[nextlevel]^[i,station]>0) then
               begin
                    Father[nextlevel,i]:=-station;
                    inc(tail);
                    OpenList^[tail].delta:=min(delta,Flow[nextlevel]^[i,station]);
                    OpenList^[tail].level:=nextlevel;
                    OpenList^[tail].station:=i;
               end;
      end;
  end;

procedure SearchWay;             { Tim kiem rong }
 var
     i,j:word;
     flag:boolean;
 begin
     repeat
       fillchar(Father,sizeof(Father),100);
       head:=1; tail:=1;
       with OpenList^[1] do { Xuat phat tu nguon }
            begin
             delta:=k; level:=0; station:=0;
            end;
       flag:=true;
       while (head<=tail) and flag do
            begin
             flag:=Open; inc(head);
             if tail>LimitDanger then { Dich hang doi de tan dung khong gian }
                  begin
                    move(OpenList^[head],OpenList^[1],(tail-head+1)*sizeof(TNode));
                    dec(tail,head-1); head:=1;
                  end;
            end;
     until flag or (k=0);        { Khong tim duong tang hoac da tim thoi gian toi uu }
 end;

begin      { Chuong trinh chinh }
     Init;
     if NoAnswer then Out(0);
     New(OpenList);
     for now:=1 to LimitTime do { Thu lan luot moi thoi diem }
       begin
           pre:=now-1;
           new(Flow[pre]);
           fillchar(Flow[pre]^,sizeof(Flow[pre]^),0); { Khoi tao luu luong }
           SearchWay;
           if k=0 then Out(now);
       end;
end.
\end{verbatim}
\endgroup

\end{document}
